\documentclass[11pt,a4paper]{article}
\usepackage[T1]{fontenc}
\usepackage[utf8]{inputenc}
\usepackage{lmodern}
\usepackage{amsmath,amssymb,amsthm,mathtools}
\usepackage{esint,graphicx}
\usepackage[a4paper,margin=22mm,headheight=14pt]{geometry}
\usepackage{microtype}
\usepackage{booktabs,longtable,array,enumitem}
\newcolumntype{L}[1]{>{\raggedright\arraybackslash}p{#1}}
\usepackage[dvipsnames]{xcolor}
\usepackage{fancyhdr}
\usepackage{hyperref}

\newtheorem{theorem}{Theorem}[section]
\newtheorem{lemma}[theorem]{Lemma}
\newtheorem{proposition}[theorem]{Proposition}
\newtheorem{corollary}[theorem]{Corollary}
\theoremstyle{remark}
\newtheorem{remark}[theorem]{Remark}
\newtheorem{definition}[theorem]{Definition}
\newtheorem{claim}[theorem]{Claim}
\setlist{nosep,leftmargin=1.5em}
\setlength{\emergencystretch}{2em}
\pagestyle{fancy}
\fancyhf{}
\fancyhead[L]{\small AIM 694}
\fancyhead[R]{\small Revised 7 November 2026}
\fancyfoot[C]{\thepage}
\renewcommand{\headrulewidth}{0.3pt}
\hypersetup{colorlinks=true,linkcolor=MidnightBlue,urlcolor=MidnightBlue,citecolor=MidnightBlue,pdftitle={AIM 694},pdfauthor={Matthew J. Colbrook; results found by AI}}
\begin{document}
\title{A translated Laguerre criterion in the positive-Gaussian case}
\author{Matthew J. Colbrook\thanks{The solutions and partial results in this document were found by AI. The supplied research and the separate repository review are AI work; human peer review and formal verification are not claimed.}\\[4pt]
\small Department of Applied Mathematics and Theoretical Physics\\
\small University of Cambridge, Cambridge CB3 0WA, United Kingdom\\
\small\href{mailto:m.colbrook@damtp.cam.ac.uk}{m.colbrook@damtp.cam.ac.uk}}
\date{7 November 2026}
\maketitle
\begin{abstract}
The strict translated Laguerre criterion is proved when the Gaussian coefficient is positive. The case of Gaussian coefficient zero remains unresolved here.
\end{abstract}
\noindent\textbf{Catalogue scope.} AIM 694; partial progress. \href{https://github.com/MColbrook/AIM/blob/1e564633494c109edfd3ec46da1366adcff925a5/problems/694-does-a-strict-laguerre-inequality-for-every-imaginary-translate-force-real-zeros.md}{Canonical target at revision \texttt{1e56463}}.


\section{Statement and argument}
\label{sec:694partial}

\textbf{Status: partial only; this problem is not counted as resolved.}
We prove the stated implication when the coefficient $a$ of the explicit
Gaussian factor in the canonical product is strictly positive.  The case
$a=0$ remains unresolved here.

For a real entire function $f$, write
\[
 U(x,y)=\operatorname{Re}f(x+iy),\qquad
 V(x,y)=\operatorname{Im}f(x+iy).
\]
Since $f$ is real on the real axis, the translated function in the
question is $f_y(x)=2U(x,y)$.  Its strict Laguerre inequality is therefore
equivalent to
\begin{equation}
 U_x(x,y)^2-U(x,y)U_{xx}(x,y)>0
 \qquad(x\in\mathbb R, y>0).
 \label{eq:694-strict}
\end{equation}

\begin{lemma}[A horizontal-decay criterion]
Suppose that a real entire function $f$ satisfies
\eqref{eq:694-strict} and, for every finite $Y>0$,
\begin{equation}
 \sup_{|y|\le Y}|f(x+iy)|\longrightarrow0
 \qquad\text{as }x\longrightarrow\pm\infty.
 \label{eq:694-decay}
\end{equation}
Then every zero of $f$ is real.
\end{lemma}

\begin{proof}
Suppose to the contrary that $f(x_0+iy_0)=0$ with $y_0>0$.
At any point of the nodal set $U=0$ in the upper half-plane,
\eqref{eq:694-strict} implies $U_x\ne0$.  The implicit-function theorem
therefore produces a local graph $x=x(y)$ through $(x_0,y_0)$ on which
$U=0$.  Continue this graph downward over its maximal interval
$(\beta,y_0]$, where $0\le\beta<y_0$.
Equivalently, it is the maximal solution of
\[
 x'(y)=-\frac{U_y(x(y),y)}{U_x(x(y),y)},\qquad x(y_0)=x_0,
\]
restricted to the nodal set.  Local uniqueness permits this continuation
as long as the graph remains in a bounded part of the upper half-plane.

Along the graph $U_x$ is continuous and never zero, hence has constant
sign.  The Cauchy--Riemann equations give
\begin{equation}
 \frac{d}{dy}V(x(y),y)
 =\frac{U_y(x(y),y)^2+U_x(x(y),y)^2}{U_x(x(y),y)}.
 \label{eq:694-monotone}
\end{equation}
Thus $V(x(y),y)$ is strictly monotone, with derivative of one fixed
sign, and its value at $y_0$ is zero.  In particular, after fixing any
$y_1\in(\beta,y_0)$, all values at $\beta<y\le y_1$ have the same sign
and have absolute value at least $|V(x(y_1),y_1)|>0$.

If $\beta>0$ and there were a bounded sequence $x(y_j)$ with
$y_j\downarrow\beta$, a subsequence would converge to $x_*$.  Continuity
would give $U(x_*,\beta)=0$ and \eqref{eq:694-strict} would give
$U_x(x_*,\beta)\ne0$.  The implicit-function theorem at this point gives
a unique graph on a whole neighborhood of $\beta$.  The original graph,
which meets that neighborhood at $(x(y_j),y_j)$, agrees with it by local
uniqueness and extends below $\beta$.  This contradicts maximality.
Consequently $|x(y)|\to\infty$ as $y\downarrow\beta$ when $\beta>0$.
Condition \eqref{eq:694-decay}, with $Y=y_0$, then gives
$V(x(y),y)\to0$, contradicting the lower bound just established.

It remains to consider $\beta=0$.  Take a sequence $y_j\downarrow0$.
If $x(y_j)$ has a bounded subsequence, pass to a further convergent
subsequence.  Since $f$ is real on the real axis, continuity gives
$V(x(y_j),y_j)\to0$ on that subsequence.  If no bounded subsequence
exists, pass to a subsequence with $|x(y_j)|\to\infty$ and use
\eqref{eq:694-decay} to reach the same conclusion.  Either alternative
contradicts \eqref{eq:694-monotone}.  There are therefore no zeros in the
upper half-plane.  Conjugation excludes zeros in the lower half-plane.
\end{proof}

\begin{proposition}
Let $f$ satisfy the hypotheses of Problem 694 and have the canonical
product
\[
 f(z)=C e^{-az^2+bz}z^m
       \prod_k\left(1-\frac z{z_k}\right)e^{z/z_k},
 \qquad
 \sum_k|z_k|^{-2}<\infty,
\]
as in \cite[Definition 1.1]{CsordasEscassut2005}.  If $a>0$, then every
zero of $f$ is real and $f$ belongs to the Laguerre--P\'olya class.
\end{proposition}

\begin{proof}
The strict inequality excludes $f\equiv0$.
Set $E_1(w)=(1-w)e^w$.  There is an absolute constant $C_0$ such that
\[
 \log|E_1(w)|\le C_0|w|^2\qquad(w\in\mathbb C).
\]
For $|w|\le1/2$ this follows by expanding
$\log(1-w)+w$; for $|w|>1/2$ use
$\log|1-w|+\operatorname{Re}w\le\log(1+|w|)+|w|$.
Fix $\varepsilon>0$ and split the canonical product into a finite head
and a tail with $C_0\sum_{k>K}|z_k|^{-2}<\varepsilon$.
The logarithm of the modulus of the finite head, including $C e^{bz}z^m$,
is at most $O(|z|+\log(2+|z|))$, uniformly in the argument of $z$.
The tail is at most $\varepsilon|z|^2$ in logarithmic modulus.
It follows that
\[
 \log\left|C e^{bz}z^m\prod_kE_1(z/z_k)\right|
 \le o(|z|^2)
\]
as an upper bound uniform in the argument.  More precisely, for every
$\varepsilon>0$ this logarithm is at most
$C_\varepsilon+\varepsilon|z|^2$.

For $z=x+iy$ with $|y|\le Y$, choose $0<\varepsilon<a$.  Then
\[
 \log|f(x+iy)|
 \le C_\varepsilon-(a-\varepsilon)x^2
             +(a+\varepsilon)Y^2,
\]
which proves \eqref{eq:694-decay}.  The preceding lemma excludes all
nonreal zeros.  The canonical-product characterization in
\cite[Definition 1.2]{CsordasEscassut2005} now gives
$f\in\mathcal{L}\mathcal{P}$.
\end{proof}

The proof uses horizontal decay supplied by $a>0$.  It does not provide
that decay when $a=0$, and it does not resolve the full conjecture.


\begingroup\small
\begin{thebibliography}{99}
\bibitem{CsordasEscassut2005} G. Csordas and A. Escassut, \emph{The Laguerre inequality and the distribution of zeros of entire functions}, Ann. Math. Blaise Pascal \textbf{12} (2005), 331--345. \url{https://doi.org/10.5802/ambp.210}.
\end{thebibliography}
\endgroup
\end{document}
